% Generated draft: requires mathematical and source review.
\documentclass[11pt]{article}
\usepackage[margin=25mm]{geometry}
\usepackage{fontspec}
\setmainfont{Latin Modern Roman}
\usepackage{amsmath,amssymb,mathtools,unicode-math}
\setmathfont{Latin Modern Math}
\usepackage{enumitem,needspace,xurl,hyperref}
\setlength{\parindent}{0pt}
\setlength{\parskip}{5pt}
\setlength{\emergencystretch}{3em}
\newcommand{\N}{\mathbb{N}}
\newcommand{\Z}{\mathbb{Z}}
\newcommand{\Q}{\mathbb{Q}}
\newcommand{\R}{\mathbb{R}}
\newcommand{\C}{\mathbb{C}}
\newcommand{\F}{\mathbb{F}}
\newcommand{\A}{\mathbb{A}}
\newcommand{\PP}{\mathbb{P}}
\newcommand{\E}{\mathbb{E}}
\newcommand{\eps}{\varepsilon}
\newcommand{\dd}{\,\mathrm{d}}
\begin{document}
\section*{UnsolvedMath ID: 30000861}
\textbf{Source problem code:} catalogue entry
\subsection{Definitions and mathematical statement}
Let $k$ be a global field with $\operatorname{char}k\ne2$, and let $(V_0,Q_0)\subset(V_1,Q_1)$ be the source's quadratic spaces of dimensions $n$ and $n+1$. With the standard nondegenerate quadratic-space convention, put $G_i=\operatorname{SO}(V_i,Q_i)$ and embed $G_0$ in $G_1$ by fixing the orthogonal complement. Let $\pi_i$ be tempered cuspidal automorphic representations of $G_i(\mathbb A_k)$, and let $\varphi_i\in\pi_i$ be nonzero cusp forms. Write $[G_i]=G_i(k)\backslash G_i(\mathbb A_k)$. The period in the question is the Petersson-type pairing \[\mathcal P(\varphi_1,\varphi_0)=\int_{[G_0]}\varphi_1(g)\overline{\varphi_0(g)}\,dg=\langle\varphi_1|_{G_0},\varphi_0\rangle,\] with the source's measure convention, and global norms are the corresponding pairings on $[G_i]$. For factorizable vectors $\varphi_i=\bigotimes_v\varphi_{i,v}$, the proposed identity is \[\frac{|\mathcal P(\varphi_1,\varphi_0)|^2}{\langle\varphi_1,\varphi_1\rangle\langle\varphi_0,\varphi_0\rangle}=2^\beta C_0\Delta_{G_1}^S P_{\pi_1,\pi_0}^S(1/2)\prod_{v\in S}\frac{I(\varphi_{1,v},\varphi_{0,v})}{\langle\varphi_{1,v},\varphi_{1,v}\rangle_v\langle\varphi_{0,v},\varphi_{0,v}\rangle_v}.\] The question asks whether an integer $\beta$ gives this formula with the predicted normalization involving the component groups (called $S$-groups in the original) of $\pi_1$ and $\pi_0$. It is not merely permission to fit an arbitrary factor separately to each pair of vectors. The extract does not give the component-group rule for $\beta$, the definition of $C_0$, the incomplete factor $\Delta_{G_1}^S$, the function $P_{\pi_1,\pi_0}^S$, the local integrals $I$, local pairings, or the admissible place set $S$. Those symbols are retained as report-defined objects, not replaced by guessed $L$-function or measure normalizations. The product notation requires decomposable vectors; the source says cusp forms without stating this restriction or a multilinear extension to general vectors.\par The statement is explicitly marked reconstructed\_unverified. Its verification note says the official report was checked but that these normalizations remain outside the extraction, so it cannot be promoted without substantial surrounding definitions. The supplied 2026-08-21 triage calls this an Ichino--Ikeda-type period formula, labels it partially solved, reports important proved cases, and says the general formula is not established in full.
\subsection{Short English statement}
Ask for the displayed normalized orthogonal-group period identity with its predicted component-group factor. The record is reconstructed and unverified because the global constants, local integrals, measures, and component-group normalization are not supplied.
\subsection{Sources}
\begin{itemize}\item Local UnsolvedMath registry, permanent ID 30000861, source problem catalogue entry, \emph{Period Formula for Special Orthogonal Automorphic Representations}, proposed year 2007; no individual proposer supplied. Statement status: reconstructed\_unverified.\item Supplied citation: \emph{Harmonische Analysis und Darstellungstheorie Topologischer Gruppen} (2008), Oberwolfach Reports, DOI 10.4171/owr/2007/49, pp.\ 2893--2942: \url{https://doi.org/10.4171/owr/2007/49}. Supplied literature label: original Oberwolfach report.\item Supplied literature label: Ichino--Ikeda period formula, \url{https://doi.org/10.1090/S0002-9947-10-05068-6}.\end{itemize}
\end{document}
